\section{Remaining Work and Upstream Path}

\subsection{Immediate Work}

\begin{table}[h]
\centering
\caption{Next work at the reporting cutoff.}
\begin{tabularx}{\linewidth}{@{}p{0.19\linewidth}p{0.18\linewidth}X@{}}
\toprule
Item & State & Notes \\
\midrule
Phase J 24h continuous & \progresspill & Daemon exists and passed a two-hour short run; formal 24-hour run remains. \\
Tier 2 / Tier 3 & \progresspill & Designs exist; templates and bootstrap flow still need implementation. \\
LTP curation & \plannedpill & Needs KEEP/SKIP list and a \UML{}-appropriate runner. \\
Snapshot readiness & \plannedpill & Unblocked after Python parity closure; should resume after Phase J. \\
Record/replay & \plannedpill & Same dependency as snapshot: architecture is planned, active work paused. \\
AVX-512 Phase B & \plannedpill & Optional follow-up for the remaining T57 vm-method failures. \\
\bottomrule
\end{tabularx}
\end{table}

\subsection{Residual Risks}

Two narrow residual flake classes are tracked as P2: an mt-mini
page-recycling class and a high-yield mt-yieldonly RIP-corruption
class. They are not current Phase J blockers and do not block
snapshot/record-replay resumption, but they should remain visible in
the report because they represent exactly the kind of concurrency
state drift the state-audit method was built to catch.

\subsection{Upstream Queue}

\begin{table}[h]
\centering
\caption{Sequenced upstream queue.}
\small
\begin{tabularx}{\linewidth}{@{}r p{0.28\linewidth}p{0.18\linewidth}p{0.18\linewidth}X@{}}
\toprule
\# & Series & Size & Subsystem & Status \\
\midrule
1 & \texttt{bpf-hygiene-v1} & 2 patches & BPF & ready \\
2 & \texttt{kmsan-arch-callback-rfc} & 1 RFC patch & mm/kmsan & ready \\
3 & \texttt{ftrace-notrace-generic-v1} & 1 patch & tracing/ftrace & prepared \\
4 & backend ops abstraction RFC & about 12 patches & arch/um & to write \\
5 & static-key hot paths & about 6 patches & arch/um & to write \\
6 & per-profile C-series & about 20 patches & arch/um plus subsystems & blocked by Series 7 ordering \\
7 & KVM backend series & about 15 patches & arch/um plus KVM & blocked by Phase J \\
\bottomrule
\end{tabularx}
\end{table}

The upstream strategy is deliberately staged. Small, independent
subsystem patches go first. The backend abstraction is the tentpole.
Static-key hot paths and profile work follow. The \KVM{} backend is
last because it is the largest technical and review burden.

\subsection{Main Message for Reviewers}

The report should not ask reviewers to accept the whole future at
once. It should make three narrower claims:
\begin{enumerate}
  \item \UML{} already has multiple backend mechanisms; making the
        backend contract explicit is a cleanup and an enabling move.
  \item static-key hooks are the right kernel-native way to make
        instrumentation runtime-flippable without charging every
        profile.
  \item the \KVMvTwo{} evidence is strong enough to justify continued
        validation and an eventual upstream RFC, but the large \KVM{}
        series should wait until Phase J is formalized.
\end{enumerate}
